\subsection{Overview}

Across the literature on agricultural subsurface (tile) drainage, authors converge on a shared meta-conclusion: despite decades of field and modelling work, the interactions among crop yield, nitrogen (N) and phosphorus (P) losses, greenhouse-gas emissions, and the field water balance remain insufficiently characterised across the full range of soils, climates, cropping systems, and drainage-management practices that occur in tile-drained agriculture. Individual practices and processes are increasingly well described, yet the field still lacks the long-term, multi-scale, and economically grounded evidence needed to design, target, and deploy cost-effective nutrient-mitigation strategies at the landscape scale. \textcite{Ma2025} distil this programme broadly, arguing that far more research is needed to understand comprehensively the impacts of conservation practices on productivity, nutrient-loss reduction, ecosystem services, and their tradeoffs in tile-drained ecosystems.

This section frames the research frontier around what remains uncertain or unstudied rather than around settled results. Several recurring lines of inquiry structure the synthesis that follows. First, there is a near-universal call for longer-duration, multi-year, multi-site monitoring to smooth climatic variability and capture the true magnitude and persistence of system responses \parencite{Randall2003,Randall2005,Waring2024,Jeong2018,Cooke2012,Kokulan2019,RANDALL1995}. Second, N management---rate, timing, form, placement, and crediting of residual and mineralisable N---remains the primary but incompletely optimised lever, together with the temporal targeting of losses \parencite{Christianson2015b,Bandura2025,Liu2021,ROSSI1991,Zhao2000}. Third, P transport in subsurface drainage is repeatedly flagged as under-studied relative to its importance, particularly for dissolved, colloidal, and bioavailable fractions and for the metrics used to predict loss \parencite{Hertzberger2019,Yuan2022,deOliveira2022,Sharpley1995,Christianson2016}. Fourth, no single best management practice (BMP) is a ``silver bullet'': bundling in-field, edge-of-field, and water-management practices and targeting them to hydrologically active critical source areas is essential \parencite{Sunohara2016,Lenhart2016,Christianson2014,Baulch2019}. Fifth, process-based models are expected to carry much of the predictive burden, yet substantial algorithmic gaps remain---especially for winter, soil carbon--nitrogen (C--N) cycling, macropore flow, and groundwater coupling \parencite{Du2017,SKAGGS1994,Dietzel2016,Feyereisen2006b}. Finally, cutting across all of these, authors stress that biophysical understanding must be coupled with economics, incentive design, policy reform, remote sensing and decision-support tools, and stakeholder engagement before mitigation can be scaled \parencite{Liu2021,Kvítek2023,David2015,Ferin2023,Khanal2020}. The subsections below organise these open questions and priorities in turn, retaining the specific quantitative targets, methods, and conditions proposed by the source studies.

\subsection{Long-term, multi-scale monitoring and experimental design}

A dominant frontier priority is the need for longer observation periods to disentangle management effects from climatic noise. \textcite{Randall2003,Randall2005} argue that long-term studies with detailed measurements are necessary to smooth year-to-year climatic variability, and stress that both the corn and soybean phases of a rotation should be evaluated when assessing N-management effects on nitrate losses. \textcite{Waring2024} similarly conclude that long-term studies are essential to capture a system's response to conservation practices and the true magnitude of that response, emphasising site-specificity, while \textcite{Cooke2012} note that studies of longer duration and many more years of yield data are required to determine drainage water management (DWM) effects on yield free of confounding factors. \textcite{Pittelkow2017} likewise stress evaluating N losses over multiple years on an annualised basis, and continuous monitoring of N application rate, timing, and method on nitrate loss and yield is advocated to give modelling a more rational setup \parencite{Jeong2018}. \textcite{RANDALL1995} add that only long-term studies across a range of climatic conditions can yield appropriate guidelines, and that plant, soil, and water measurements must be gathered together to uncover interactions among crop production, climate, and environmental impact.

Several studies extend this call to non-growing-season and multi-year hydrologic coverage. \textcite{VanEsbroeck2016} recommend collecting field data during the non-growing season when estimating P loads and BMP efficacy, and call for data across a wider range of soil textures, management practices, and hydrologic years---including wetter summers---before drawing conclusions on optimum P-application timing. \textcite{vanderSalm2012} argue that data spanning at least a normal, a wet, and a dry year are necessary to understand losses and transport routes, and \textcite{Ulén2007} note that as long a record as possible is needed for trend estimation because water-quality change may lag the P-loss change driven by agricultural measures. Long-term or continued monitoring is repeatedly invoked to evaluate P leaching and nutrient removal of infiltration systems \parencite{Duan2015}, to determine whether enhanced effluent-driven P losses persist \parencite{Monaghan2004}, to confirm bioreactor and saturated-buffer effectiveness across soils, landscapes, and weather years \parencite{Jaynes2019}, and to obtain definitive evidence of suspended-solids and orthophosphate reductions under cover cropping \parencite{SalcedaGonzalez2025}. Long-term designs are also invoked to detect slow or subtle change---in soil organic carbon recovery under rotation shifts \parencite{Kumar2014}, in discharge-induced organic-matter export \parencite{Manninen2023}, and in whether drainage-induced soil evolution continues or stabilises over decades \parencite{Frison2009}. Several teams describe active or planned long-term programmes in rice--wheat and rice--oil rotations \parencite{Qi2025b}, in subsurface pipe and seepage drainage across diverse environments \parencite{Liu2025b,Liu2025c}, and note that four years of data are insufficient to resolve residue-cover effects on pollutant losses and yield \parencite{Thoma2005,Acharya2019b}.

Experimental-design refinements are proposed alongside longer records. \textcite{Schilling2013} recommend establishing a suitable correlation among paired basins across a range of climates, restricting paired sampling to the March--July high-flow period, delaying post-BMP monitoring to enhance the likelihood of detecting change, and using BMPs capable of achieving at least the minimum detectable change. \textcite{Sunohara2016} argue that BMP impacts on downstream receptors should be assessed simultaneously at multiple scales, from field to watershed, and \textcite{Irvine2019} call for within-field and within-plot studies to elucidate mechanistic processes of nutrient export and for additional field-scale studies within high-export sub-watersheds. Further experiments under wider conditions are sought to test effective-yield formulas and models \parencite{Xin2016c}, refined $^{15}$N loss observation over larger, more representative trial areas is advocated \parencite{Xiong2025}, and mesocosms are seen as an opportunity to quantify nutrient-loss responses to management change \parencite{Loper2024}. Underlying these appeals is the imperative that longer records are only useful if sampling is standardised: soil biological and soil-health indicators should be sampled at consistent times under similar moisture and temperature conditions \parencite{Sherbine2023,Li2023b}, and meta-analytic comparability requires complete hydrologic description of each plot, standardised analytical methods, adequate monitoring periods, and full reporting of concentrations with descriptive statistics \parencite{Ghirardini2019}.

\subsection{Nitrogen management: rates, timing, form, crediting, and loss timing}

Refining fertiliser N management is the most consistently identified priority. \textcite{ROSSI1991} state the urgency plainly: management directed at substantial nitrate reductions is needed, with careful evaluation of the effects of fertiliser timing, rate, form, and placement. \textcite{Zhao2000} argue efforts should concentrate on strategies that increase N-use efficiency in corn---optimal rate, timing, and form---while \textcite{Randall2001} call to apply the correct rate at the optimum time, credit previous legumes and manure, and improve manure management through known nutrient content, uniform spreading, and timely incorporation. At watershed scale, \textcite{Dattamudi2020} conclude that N-application schemes need reevaluation, suggesting lower corn rates, precision agriculture, use of soil N with high-yielding varieties, and split, properly timed applications; \textcite{Elmi2002} likewise see reductions from matching rate and timing to crop uptake. Optimising N rates is expected to continue receiving primary research and regulatory focus, and accurate 4R implementation will require site-specific knowledge to balance agronomic and environmental goals \parencite{Christianson2015b}.

\subsubsection{Next-generation rate tools and decision support}

\textcite{Bandura2025} propose concrete standards for the next generation of N-rate tools: they should be developed and refined for various application timings using the economically optimum N rate (EONR) as a benchmark, judged adequate if they predict N need within 10~kg~N~ha$^{-1}$ of EONR, and environmentally unacceptable if they exceed EONR at an N:corn price ratio of 5.6 by more than 30~kg~N~ha$^{-1}$. \textcite{Liu2021} call for more accurate recommendations that weigh environmental and economic benefits, effective local nutrient plans, case-by-case cost analyses, and coordinated agency action (e.g.\ USDA-NRCS with USEPA) to promote 4R adoption. \textcite{Banger2018} advocate additional experiments across wider management and soil types to identify potential yield tradeoffs from rate reductions, noting that if yield-loss risks appear, emphasis should shift toward cover crops and edge-of-field practices to meet load-reduction targets. Comprehensive tools integrating reliable weather forecasts, historical management, local field and soil parameters, and crop-variety information are proposed to estimate optimum rates, alongside regional long-term agronomic-optimum-N-rate (AONR) means and study of the planting-date/efficiency/AONR relationship \parencite{Maas2022}, and precision application based on soil-N availability \parencite{Gupta2023}. \textcite{Dinnes2002} prioritise better crediting of non-fertiliser N sources, improved in-season soil-N monitoring for split applications, and reliable variable-rate models incorporating remote sensing and hyperspectral approaches, while \textcite{Cassman2025} propose low-cost proxy variables---yield without applied N as a field-scale estimate of soil N-supply capacity, and partial factor productivity of applied N and rainfall---for monitoring soil N and water supply.

Local-scale testing and decision approaches require refinement. Late-spring nitrate-test (LSNT) programmes are framed as an initial start to be replaced or refined as new research emerges, with high-clearance equipment to reduce delayed-sidedress risk and cropping diversity or edge-of-field denitrification needed to reach drinking-water goals \parencite{Jaynes2004}; grain-quality and grain-protein monitoring plus risk insurance are recommended for LSNT adopters \parencite{Karlen2005}; and weather-based model refinement into simple tools, with improved understanding of year-to-year soil mineral-N variability, is proposed to reduce soil-testing needs and increase side-dressing adoption \parencite{Malone2010}. \textcite{Jaynes2013} call for research to pinpoint where between the V12 and V16 growth stages sidedress N ceases to be beneficial, and \textcite{Jaynes2015} note sidedressing can be combined with soil tests or crop N-deficiency sensing for real-time adjustment. Broader gaps include comprehensive volatilisation-loss studies, targeted research from data-poor regions, and provision of ancillary climate, yield, and soil data to support predictive modelling and policy \parencite{Chatterjee2020b}.

\subsubsection{Timing, form, and application method}

Timing and method of application remain recurring open questions. \textcite{Liang2018} propose combining timing with rate and note the value of models such as WHCNS for evaluating short- and long-term effects. \textcite{Malone2007b} highlight the need to quantify differences between split and single spring applications and to check possible over-prediction of winter nitrification and N loss for fall manure on soybean, while \textcite{HernandezRamirez2011} note it remains unknown whether reduced-rate fall manure could match sidedress or preplant applications environmentally, pointing to nitrification inhibitors, cover crops, more complex rotations, and multifunctional landscapes as additional avenues. \textcite{Elmi2005} stress quantifying all major N loss pathways simultaneously and obtaining better information on manure versus mineral systems. Product- and method-level questions are increasingly prominent: enhanced-efficiency fertiliser technology---spring-applied anhydrous ammonia with a nitrification inhibitor---should be tested in conjunction with enhanced DWM \parencite{Kaur2024c}; NH$_4^+$-N management via N stabilisers, and cover-crop effects on field- and watershed-scale NH$_4^+$-N dynamics, deserve exploration \parencite{Vincent2025}; and fertigation of silage corn requires study of whether drainage nitrate can be reduced without yield loss, given tradeoffs of crop injury, disease, and pathogen transport relative to slurry injection \parencite{Gamble2018}. Evening out manure application within high-livestock catchments, and reducing spring and autumn soil mineral N in cereal areas, are further priorities \parencite{Bechmann2014}. \textcite{Haghnazari2020} call for comparing the economic and environmental consequences of N use under global-warming projections.

\subsubsection{Crediting residual, mineralisable, and legacy nitrogen}

A recurring, more specific gap is the crediting of N supplied by soil organic matter, prior legumes, manure, and land-use legacies when setting rates. \textcite{Bakhsh2007,Bakhsh2005} call for assessing proper N credit from mineralisation of the soil organic-matter pool and soybean residue, and the storage of mineralisable N in the root zone, before applications. \textcite{Randall2001} extend this to improved soil-N testing methods to estimate mineralisable and carryover N, especially after dry years, legumes, or manure---echoed by a recommended spring soil-N test following dry years to credit residual nitrate \parencite{RANDALL1995} and by measuring mineral N to evaluate leaching potential before sowing \parencite{Conrad2009}. In systems transitioning from perennial cover, \textcite{Huggins2001} propose second-year N credits for mineralisation following grass conversion, and \textcite{ROBBINS1980} counsel care in setting rates and timing on land plowed out of alfalfa. For irrigated vulnerable soils, \textcite{Derby2009} recommend reducing N in early irrigation years, setting reasonable yield goals, crediting fall residual nitrate, and pursuing field-scale studies spanning soil, vadose-zone, and groundwater. The role of \emph{legacy} N is itself flagged as underexplored: a clearer definition of legacy nutrients is urgently needed to apportion losses \parencite{Williams2022}, and the coupled biogeochemical and hydrological processes regulating legacy-N formation and its impact on nitrate export in tile-drained agroecosystems require dedicated study \parencite{Hu2024}. Better quantification of fertiliser rates and timing, and reconsideration of fall application under warmer winters, would sharpen estimates of management influence on nitrate transport \parencite{Gentry2014}, and further quantification of soil N in growing and non-growing seasons is advised \parencite{Bosompemaa2021}. Emerging tooling---resin-membrane strips to characterise soil-N cycling \parencite{Masri2024}---and better predictions of long-term flow-weighted N concentrations for benchmark soils under typical rotations \parencite{Skaggs2012} round out this agenda.

\subsubsection{The temporal window of loss}

A distinctive frontier is the temporal concentration of losses, particularly in early spring, and the need for management targeted to that window. \textcite{Bakhsh2007} state that strategies must be developed to reduce early-spring nitrate leaching, and \textcite{Bakhsh2010} extend this to the March--June period, adding that the temporal distribution of growing-season rainfall needs further attention. Assessment of soil nitrate status, tile flow, and nitrate loss during April, May, and June warrants further study \parencite{Nance2007}, and a reliable winter-hardy cover crop ahead of corn, together with a shift of fall N to spring, may be the best way to decrease loss in intensively tile-drained watersheds \parencite{Gentry2024}. For outlet management, \textcite{Skaggs1995} recommend selectively managing drainage outlets to reduce peak nitrate losses, especially in years following drought-reduced yields, and both \textcite{Elmi2004,Elmi2002} note that keeping controlled drainage operational during early spring and late fall---when losses are most severe and drainage is not needed for production---could achieve further reductions.

\subsection{Phosphorus: metrics, mechanisms, transport, and prediction}

Phosphorus in subsurface drainage is repeatedly identified as under-researched relative to N. \textcite{Hertzberger2019} explicitly call for greater research effort on P transport and better understanding of nutrient-related water-quality issues, while \textcite{Yuan2022} note more research is needed on conservation-practice effectiveness in reducing dissolved-P loss in subsurface flow and on the causes of annual variability in reduction rates. \textcite{Pease2018} identify the pedologic factors influencing P transport and the effect of management on P loss as key needs across the Eastern Corn Belt, and \textcite{deOliveira2022} recommend refining P-loss frameworks through additional within-field studies, continued data collection on soil testing, site history, and agronomic management, and better understanding of the ``where'' and ``when'' of drainage P loss to sharpen conservation-practice timing and placement.

\subsubsection{Beyond agronomic soil-test P: saturation and sorption metrics}

A consistent theme is that agronomic soil-test P (STP) is insufficient for environmental risk. \textcite{Sharpley1995} argue that agronomic tests may not reliably assess eutrophication potential and that environmental tests estimating bioavailable P and sorption capacity should be applied to high-risk soils. \textcite{Hooda1999} propose that degree of P saturation (DPS) may better estimate loss potential than extractable P such as Olsen-P, and that a full accounting of all P output forms is needed; \textcite{Nguyen2002} call for establishing critical DPS values calibrated across sediment P tests and equilibrium P concentrations. \textcite{Osterholz2020} recommend assessing additional fields in the 100--400~mg~kg$^{-1}$ STP range and examining sorption saturation, macropore formation, and mineralogy, while P-adsorption capacity should be considered before organic soils are developed for cropping \parencite{MILLER1979}. \textcite{McDowell2001} propose that P isotope-exchange kinetics (IEK) has the potential to explain and eventually predict P loss to overland flow and subsurface drainage across a wide range of STP concentrations.

\subsubsection{Desorption, colloidal and particulate transport}

Multiple authors identify the sorption/desorption interface and particulate pathways as poorly resolved. \textcite{Mumbi2025} regard quantifying when and how varying water sources and fluxes interact with soil to desorb and adsorb P as critical to advancing prediction of subsurface DRP transport, and \textcite{Flores2021} note the need to understand mechanisms controlling particulate settling and P desorption. \textcite{Uusitalo2003} argue that directing runoff into the subsurface pathway and establishing field-margin buffers are only first steps, that mitigating bioavailable particulate P requires reducing particle mobilisation across the entire field surface, and that controlling diffuse DRP is a longer-term task than reducing erosion---with knowledge of the availability of eroded P to algae still limited \parencite{Ulén2007}. \textcite{Wilson2019} highlight the need to quantify the contribution of colloidal P to DRP measurements and to determine whether erosion-mitigating BMPs or nutrient management better control it, alongside clarifying stratification, preferential flow, and the effectiveness of subsurface P application. \textcite{Norgaard2013} recommend including both dissolved and colloid-associated leaching parameters in field-scale risk assessment. Sediment P release under prolonged flooding and alternating waterlogged--drying regimes \parencite{Nguyen2002}, P fractionation after leaching under polyacrylamide \parencite{Jiang2010}, and combining intact columns, plots, and catchments to examine subsurface loss of particulate matter and P in relation to macroporosity \parencite{Laubel1999} are all advanced as priorities.

\subsubsection{Management effects, events, ratios, and legacy P}

\textcite{Uusitalo2018} note that assessing the pace of P stratification requires more data, suggest simulation tools such as the Annual P Loss Estimator for evaluating mixing intervals, and leave open whether no-till adoption as a eutrophication-control measure is justified. \textcite{Ni2020} call for quantifying dominant mechanisms of agricultural conservation practices under complex scenarios and for studies of optimal soil-test-P levels under changing climate to address STP--precipitation interactions, while \textcite{Wang2018} ask whether pre- or immediate post-application tillage decreases DRP losses through tile drainage and whether results apply to manure injection. Manure-related risk---the immediate P loss after a single application---needs fuller understanding \parencite{Madison2014}. \textcite{Brendel2019} recommend designing BMPs that reduce P losses during frequent small events, understanding relationships between flow, precipitation, pathway, and P fraction, and using these to target implementation, and \textcite{Schilling2017b} call for refining organic-P/total-P ratios, quantifying wastewater-treatment-plant contributions, and identifying sources of variability such as manure and animal density. \textcite{Jarvie2017} ask whether changing climate contributes to increased soluble reactive P delivery, note it is unknown how much of a 40\% reduction target could be achieved by 4R practices alone, and warn that drawing down legacy P and disconnecting tile-flow P loads at times of maximum receiving-water sensitivity may be increasingly costly and intrusive. Cost-effective mitigation hierarchies are proposed, prioritising lower soil Olsen P and targeting alum to critical source areas \parencite{McDowell2015}, and P removal via regular grass harvesting under continuous-grass tile-drained fields \parencite{Zhang2015b}.

\subsubsection{Multi-species reporting, tillage, and databases}

Fragmented reporting hampers synthesis. \textcite{Christianson2016} call for more field-scale studies documenting drainage P losses across cropping, nutrient, soil, and design factors, including multiple P species (DRP, particulate, total), and for further developing the Drain Load database. \textcite{Macrae2019} seek understanding of hydroclimatic drivers of event responses and how tillage or P placement affects concentration relationships, and \textcite{Tanga2025} recommend monitoring both particulate and dissolved forms while selecting site-specific stoplog depths from slope, soil, and climate. \textcite{Zhang2017} add that impacts of modified no-till on DRP need investigation and that new soil-P testing and optimised application timing, rates, and methods are required under no-till. Modelling refinements proposed include adjustable manure-P mineralisation and decomposition parameters and further tillage-effect investigation \parencite{Pan2023}, continued P-loss model improvement and representation of crop stress prior to maturity \parencite{Wang2019}, and applicability testing of in-situ optical turbidimeters for particulate matter in drainage water \parencite{Grant1996}. \textcite{Daigh2015b} note that reactive-P concentrations and yields in bioenergy systems need further evaluation where P losses are high, and \textcite{Spiese2023} call for novel BMPs addressing both nitrate and DRP simultaneously.

\subsubsection{Phosphorus in organic (Histosol/peat) soils}

Phosphorus in cultivated organic soils is singled out as poorly understood. \textcite{Grenon2022} call for mechanistic studies of the pathways, hot spots, and hot moments of P transfer from agricultural organic soils, including residual-fraction dynamics, long-term soil--water study, and sorption dynamics for accurate balances, and note that more frequent post-fertilisation sampling and isotope tracers would help. \textcite{Grenon2021} propose a concrete agenda: assessing P geochemistry in Histosols using isotope tracing ($^{32}$P/$^{33}$P and oxygen isotopes in phosphates), XANES, and $^{31}$P~NMR alongside sequential fractionation; assessing microbial turnover; long-term DWM studies to reduce total-P loads; a new model accounting for geo-mechanical instability; and P-transport indices. Nutrient leaching from abandoned peat fields is likewise poorly studied, calling for research into ditch blockage, energy-willow, or forestation, and into the true nature of H$_2$SO$_4$-P via P K-edge XANES \parencite{Höyhtyä2025}.

\subsection{Cropping systems: cover crops, perennials, rotations, and diversification}

\subsubsection{Cover crops: agronomic, environmental, and modelling frontiers}

Cover crops---especially winter rye---attract the largest concentration of open questions. \textcite{Kladivko2014} call directly for further research, development, and technical assistance to overcome practical obstacles to adoption, for additional field experiments on multiple soil types across the Midwest, and for using new data to calibrate and validate RZWQM. \textcite{Birru2023} and \textcite{Zhao2016} stress that existing evidence is drawn from too few, geographically clustered sites and call for extending scope across more species, soils, climates, and management. On the agronomic side, priorities include tools to promote rye growth in colder regions and long-term no-till study \parencite{Waring2020}, strategies to establish rye under temperature, precipitation, and growing-degree-day extremes \parencite{Bawa2023}, evaluation of more cultivars and improved establishment to reduce year-to-year variability and yield impacts \parencite{Kaspar2007}, cultural and plant-breeding research to improve compatibility and cold hardiness \parencite{Dinnes2002,Liu2019}, and investigation of whether observed yield drag reflects poor planter performance through residue \parencite{Dougherty2020}. Kura clover as living mulch with corn requires better agronomic management to avoid yield loss \parencite{Qi2011b,Qi2011}.

On environmental mechanisms, authors call for research on the release of scavenged N in spring and its synchronisation with corn and soybean growth stages \parencite{Lacey2015}, producer-adoptable methods to determine cover-crop N content and adjustable seeding rates \parencite{Roth2018}, the combined effects of cover crops and manure rate on subsurface DRP \parencite{Askar2023}, and optimisation of the trade-offs between cover-crop services and disservices for the water balance \parencite{Meyer2019}. \textcite{Liu2019} more broadly call for coordinated common experiments across geophysical conditions to quantify winter-cover P contributions, identify reliable analytes, and predict future-climate influences. Longer-term evidence remains scarce: additional long-term studies of winter cover-crop effectiveness \parencite{Kaspar2012}, longer study of rye influence on soil-health parameters \parencite{Bawa2023}, and understanding of the long-term consequences and biotic--abiotic mechanisms by which rye affects yield and environmental performance \parencite{MartinezFeria2016} are all sought. Optimising cereal-rye management \parencite{deOliveira2022}, growing cover crops longer in spring to prevent N$_2$O increases \parencite{PrezaFontes2023}, and identifying practices that target individual nutrient pathways because over-winter effects differ for N versus P \parencite{McNeill2026} are further priorities, while \textcite{Krueger2013} recommend conservation practices beyond rye because subsurface nitrate remained high and rye had no effect on DRP or drainage volume.

Cover-crop \emph{modelling} is a distinct frontier. Needed advances include validation of the rye growth model and better plant-N-uptake algorithms \parencite{Feyereisen2006}, thorough multi-year testing to fix winter cover-crop parameters and more realistic N-uptake routines \parencite{Malone2007c}, improved simulation of rye germination and establishment \parencite{Malone2014b}, a dedicated winter-rye crop module and careful calibration of soil-organic-matter and drained-upper-limit parameters \parencite{Chatterjee2025}, correction of APSIM weaknesses in cover-crop effects on following corn, frozen-soil temperatures, and soybean drainage \parencite{Dietzel2016}, a chemical-termination routine \parencite{Gupta2022}, and improved drainage and mineralisation/immobilisation modules with attention to extreme events \parencite{Gupta2023b}. Additional species such as radish and legumes are proposed for hydrology and water-quality simulation \parencite{Singh2023}, and lysimeter-paired field validation of rye effects on soil-water dynamics and nitrate leaching is underway \parencite{Qi2010b}.

\subsubsection{Perennials, bioenergy, and rotational diversification}

Perennial and land-use-change options feature prominently. Poorly drained depressions are proposed as candidates for perennial bioenergy crops such as \textit{Miscanthus $\times$ giganteus} or wetland restoration \parencite{Huang2023}, and keeping soil covered year-round with a growing crop such as perennial ley is judged the most effective way to mitigate N losses, while the P focus should be on decreasing surface runoff \parencite{Norberg2022}. Additional data on alfalfa's effect on subsurface discharge during large events across diverse geographies are sought \parencite{Arrueta2025}, and a clear causal link between perennial-crop implementation and nutrient loads across soils is called for, alongside policies rewarding environmental services \parencite{Arrueta2022}. Robust observed water-quality datasets for perennialised systems \parencite{Wilson2023}, large-scale assessment of bioenergy-crop scenarios with cost--benefit analysis to alleviate Gulf of Mexico hypoxia \parencite{Guo2018}, and research on the value of and payments for externalities such as P, sediment, greenhouse-gas, and biodiversity benefits of switchgrass \parencite{Woodbury2018} extend this agenda. Rotation and diversification questions recur: constraining crop-rotation impacts on yield-scaled N loss and accounting for losses on a rotation rather than annual basis \parencite{Zhao2016}, the paucity of research on desirable rotation crops in two-crop Midwestern systems \parencite{Pasley2021}, alternative perennial-containing systems and partial conversion of cropland to pasture or wetland to attain nitrate TMDLs \parencite{Randall2001,Gowda2007,Nangia2008}, and explicit evaluation of controlled drainage, rotation, cover crops, and N management together \parencite{Ma2007b}. Novel options such as intercropping triticale with red clover and developing winter triticale \parencite{Nance2007}, camelina--soybean relay- and double-cropping \parencite{Emmett2022}, fertilised/harvested rye for revenue and net energy \parencite{Malone2018,Malone2023}, and a diverse cropping system documented against the conventional two-year rotation \parencite{Gentry2025} are all advanced for study. \textcite{Dinnes2002} call broadly for perennial cash-cropping systems that maintain returns while reducing leaching, and \textcite{Jaynes2001} conclude that dramatic management change on artificially drained soils will be required before corn production can be considered sustainable.

\subsection{Drainage water management, water-table control, and system design}

Controlled drainage (CD), shallow drainage, subirrigation, and other water-table management approaches are seen as promising but requiring optimisation, verification across climates, and clearer mechanistic partitioning. \textcite{Jaynes2012} are unaware of research optimising DWM operation for the soils, climate, and crops of central Iowa, and call for identifying locations where more field area falls under control-structure influence and for cost-sharing incentives. \textcite{Fang2012} stress that careful management of gate timing and depth is essential to optimise reduced N loss, and \textcite{Wesström2001} argue that water balances for conventional and controlled plots must be established before differing outflow patterns can be separated. Whether managed drainage can consistently produce greater yield than free drainage---perhaps concentrated in higher-yield-potential years on claypan soils---remains unresolved \parencite{Nash2015,Nelson2011}, and continued study in normal-precipitation years may clarify whether CD offers water-conservation advantages and whether shallow drainage reduces nitrate \parencite{Smith2011b}. CD adoption may become more relevant under wetter, warmer climates, requiring practical control-structure management during wet springs \parencite{Madramootoo2024}, and the collective impacts of automated DWM on the coupled timing of flows and composition of exported water need study before its full BMP efficacy can be assessed \parencite{Miller2022}. Optimum dormant-season management depth needs more experience \parencite{Adeuya2012}, better on-site monitoring, long calibration periods, and anti-seep collars are recommended \parencite{Gunn2015}, and economic and agronomic questions---yield effects via SEW$_{30}$, planting-date impacts, and the economic risk of reducing nitrate loss with DWM---remain open \parencite{Schott2017b}.

Optimising drainage design for the joint goals of yield and water quality is a persistent problem. \textcite{Kladivko2021} argue that more field studies coupled with modelling are urgently needed to develop site-specific recommendations and optimise design for both yield and water quality, and \textcite{Zhao2016} suggest research into alternative drain spacing balancing production against decreased losses. Field experiments characterising how shallow and controlled drainage affect the complete water balance, including soil-moisture-based flexible control \parencite{Singh2006b,Singh2007b}, low-cost retrofittable water-table-control designs \parencite{Dinnes2002}, and subirrigation extended to diverse regions with in-depth economic assessment \parencite{Gunn2019,Baker2004} are all sought. Cold-climate and peatland settings raise specific needs: optimising winter subsurface drainage and water-table management as temperatures warm \parencite{Li2025}, evaluating subsurface irrigation to maintain shallow water tables and mapping cultivated peatlands \parencite{Läpikivi2025}, and cautions that severe frost may irreversibly affect ploughed-soil structure \parencite{STEINHARDT1974}. Drainage \emph{materials and installation} also feature: long-term investigation of knitted-sock envelope clogging \parencite{Ghane2022}, straw and crop-residue drainage materials and integrated excavation-and-burial machinery \parencite{Qi2025b,Lu2018}, DRAINMOD-based evaluation of design impacts on production, salinity, and waterlogging \parencite{Ellafi2021}, countermeasures such as combining drainage with raised-bed planting \parencite{Takeshima2023}, and perforation-characteristic and standardisation studies for drain materials \parencite{Zhang2025b,ASGHAR1995}. Nutrient-trading programmes are identified as needed to increase shallow-drain adoption \parencite{Ghane2021b}, and in paddy and rice systems new water-management options may reshape agronomy and require new infrastructure and precision integration \parencite{DarziNaftchali2014,Jung2024,Hashemi2020}.

\subsection{Edge-of-field and in-field treatment technologies}

Engineered nitrate- and P-removal structures---bioreactors, saturated riparian buffers (SRBs), constructed wetlands, and filter-based drainage---constitute a frontier where design optimisation and long-term field performance remain open. \textcite{McEachran2020} state that a key challenge is to predict nitrate removal so that more effective SRBs can be designed, outlining methodological advances (unsteady governing equations, average linear velocity, numerical nitrate transport, and optimal-width approximation). \textcite{Jaynes2014} propose that replicating SRB results across landscapes could prevent substantial nitrate from entering midwestern streams and that biomass harvesting might maintain long-term viability, while \textcite{Abdalaal2023} recommend more sophisticated removal and hydrology components, coupling design with economic or objective-function targets, and determining maximum allowable buffer length before optimising width. Confirming that SRB N removal increases as vegetation matures, and quantifying gains from individual buffer components with monitoring wells and backflow valves, are further needs \parencite{Jaynes2019,Singh2024b}. For bioreactors and combined installations, \textcite{Husk2017} call for study of CD with and without a denitrification bioreactor to quantify each component, note that treating greater volumes at lower efficiency could improve mass removal, and warn that clogging potential must be evaluated. Refining denitrification-bioreactor design---reactor size, packing materials, equilibrium thresholds, elimination of preferential flow, and pH maintenance---is emphasised \parencite{Aide2016}, as are vegetation and organic matrices in reactive riparian zones \parencite{Vallée2014}, wetland restoration or denitrifying ponds to treat drainage before it enters ditches \parencite{Randall2001}, and a broad agenda for constructed treatment systems spanning driving factors, flow--treatment relationships, design-parameter effects, long-term field tests, and winter-efficiency enhancement \parencite{Mendes2021}. For filter-based drainage, nitrification and denitrification in soil and filter cannot be neglected over long operation, and long-term experiments with different filter materials are needed \parencite{Tao2022}.

\subsection{Bundling, sequencing, spatial targeting, and critical source areas}

A pervasive conclusion is that no single BMP suffices and that combinations must be studied, sequenced, and optimised. \textcite{Sunohara2016} state plainly that no one BMP is a silver bullet, and \textcite{Lenhart2016} argue that multiple BMPs in series are needed, that in-field/edge-of-field combinations show great promise but have been little studied, and that wetlands and water-retention BMPs must be developed for sensitive downstream waters. Attaining water-quality goals will require a combination of practices implemented at a high level in a targeted manner \parencite{Christianson2014}. Combining conservation tillage with cover crops, pairing in-field with edge-of-field conservation, and incorporating legacy effects \parencite{Makhtoumi2025}, examining synergistic effects and formalising optimisation studies \parencite{Kreig2019}, and better understanding of practice interaction and the N-removal capacity of various practices with improved mapping of tile-drained land \parencite{McLellan2015} are all advanced. \textcite{Baulch2019} lay out a broad cold-region programme---comparing surface, tile, and undrained fields; studying CD, channel design, and maintenance through wet--dry cycles; assessing tillage and perennial-forage conversion; quantifying retention in reservoirs, wetlands, and buffers; and developing snowmelt-suited tile-drainage mitigation. Wider application of bioreactors, P reactors, CD, and low-impact development, potentially alongside internal-loading management \parencite{Husk2024,Husk2017}, robust cost-effective BMPs to support cover crops and mitigate warming \parencite{Gupta2023}, and intercrops and hydrology-altering management for P-transport reduction \parencite{HernandezRamirez2011,GAYNOR1995b} extend the bundling agenda.

Closely related is targeting practices to the parts of the landscape that matter most. \textcite{Melland2008} recommend identifying hydrologically active areas as mitigation priorities and avoiding high stocking and inputs in waterlogged runoff-generating zones. Whether elevated loads originate from numerous low-intensity or a few key hot spots \parencite{Irvine2019}, whether targeted or regionally saturated cover-crop implementation is required for consistent nitrate and soluble reactive P reductions \parencite{Speir2022}, and the delineation of zones transporting higher nitrate leaching that require site-specific management \parencite{Bakhsh2005b} are all open questions. Targeting critical source areas and hydrologically connected hotspots is repeatedly stressed \parencite{Baulch2019,McDowell2015,Brendel2019}, prioritisation is expected to improve return on conservation investment \parencite{Makhtoumi2025}, and automating spatial-targeting principles into an easily applicable GIS tool with whole-basin surveys and cost--benefit analysis is prioritised \parencite{Kvítek2023}. Combining GIS with hydrological modelling and saturation risk to assess DRP and nitrate loss, and finding the nitrate threshold across P saturation, texture, and temperature, are further scaling priorities \parencite{Smith2023}, alongside better understanding of the coupled effect of controlling factors and drainage density on stream solute dynamics \parencite{Kennedy2012} and identification of critical precipitation periods rather than annual values \parencite{Salo2006}.

\subsection{Process-based modelling: outstanding algorithmic and evaluation gaps}

Because field trials cannot cover all combinations of practice, soil, and climate, process-based models are expected to fill the gap, yet the literature documents extensive, still-unresolved modelling needs. \textcite{Du2017} anticipate models will be relied upon extensively as concern over nitrate export increases, and \textcite{SKAGGS1994} caution that much remains to be learned about the mechanisms governing pollutant losses, that process understanding must improve before quantification is possible, and that better datasets are needed to test models across wide ranges of soil, crop, and climate. \textcite{Smith2008} argue that a validated process-based rather than empirical model should be applied to capture complex variable interactions.

\subsubsection{Soil carbon--nitrogen cycling, plant uptake, and crop growth}

Many authors identify soil N-process representation as the weakest link. Priorities include improving nitrification, denitrification, fermentation, and mineralisation within an improved hydrology framework \parencite{Smith2020}, better representing crop N uptake, soil N cycling, and C--N pool transformations to reduce errors in flow-weighted nitrate predictions \parencite{Fang2012}, and improving soil N processes, soybean N fixation, and maize-yield simulation in SWAT \parencite{Hu2007}. RZWQM improvement agendas recur throughout: accounting for subsurface lateral flow below tiles, simulating rather than supplying soybean N fixation, improving plant N uptake, and reducing winter nitrification so more NH$_4$-N remains at planting \parencite{Ma2007}; better simulation of management effects on yield and mineralisation and extension to additional tile-drained sites \parencite{Ma2007b}; improved crop responses to N-application timing and rigorous evaluation of each C/N process before microbial-growth roles can be understood \parencite{Ma2008}; and a wide-ranging programme spanning variable cracking, P dynamics for manure, mobile organic-matter movement, new-crop parameterisation, standardised calibration, and GIS/economic linkage \parencite{Ma2001}. Modifying SWAT nutrient-cycling algorithms for immobilisation under corn-stover removal \parencite{Malone2019}, revising winter processes and organic-C mineralisation in the PNM model \parencite{Marjerison2016}, a more robust approach to grass biomass and N uptake in drained grassland \parencite{Bechtold2007}, refining N-cycling for N$_2$ fixation and plant uptake to correct soybean over-prediction \parencite{Bakhsh2004}, and a more physically based tile-nitrate model using ``soft data'' and a passive humus pool \parencite{Qi2025} complete this cluster. Systematic data collection on water and solute fluxes, greenhouse-gas emissions, and root and biomass development is repeatedly identified as the foundation for model validation and refinement \parencite{Ma2025}.

\subsubsection{Winter, frozen-soil, and snow processes}

Cold-season representation is singled out as inadequate. Better algorithms for soil freezing/thawing and spring-thaw drainage, and more detailed mechanistic plant-N-uptake and snow--freeze--thaw representation, are needed \parencite{Feyereisen2006b}. DRAINMOD-N needs improvement for freezing and thawing seasons and additional inputs such as tillage and intercropping \parencite{Yang2007}, and winter processes in seasonally frozen regions warrant investigation \parencite{Saadat2020,Jiang2020}. An impedance factor for frozen-soil hydraulic conductivity, ice-content measurements, and a calibratable snow--rain temperature threshold are proposed for the RZ-SHAW model \parencite{Li2024b}, and continued development and testing of RZWQM for winter conditions---including snow algorithms and improved water-retention parameterisation---is recommended \parencite{Ahmed2007b,Ahmed2007,Yang2021}. Swelling and shrinkage during infiltration into frozen soil \parencite{Roy2020} and improved winter-crop parameterisation with macropore infiltration evaluation \parencite{Motarjemi2021} round out this agenda.

\subsubsection{Macropores, preferential flow, and drain configuration}

\textcite{Smith2020} recommend improving runoff simulation under surface crusting, clay cracking, preferential flow, snow dynamics, and freeze--thaw, and \textcite{Stamm2002} argue that more attention should be given to actual drain configuration to complete the picture of solute transport in drained soils. SIDRA-RU upgrades---a new nitrate-transport model, sampling-interval sensitivity analysis, crop-rotation coupling, and time-variant parameter estimation \parencite{Chelil2022b,Jeantet2021}---and runoff models that include temporal variability in hydrological-pathway activation for cracking clay soils \parencite{Kokulan2019} are proposed. \textcite{Boivin2006} call for focus on hydraulic-conductivity and shrink-swell parameters and their seasonal time-dependence, direct in-field measurement methods, and delineation of seasonal preferential-flow processes.

\subsubsection{Drain-design decision support, coupling, and validation}

\textcite{Moriasi2013} propose using the Revised SWAT model to determine optimum drain depth and spacing but stress that validation of the Hooghoudt and Kirkham equations is needed before use as a decision-support tool. Validating nitrate-load methodologies against DRAINMOD-NII or RZWQM and in additional watersheds \parencite{Ale2012}, multi-year calibration for free-drainage treatments \parencite{Ng2000}, additional testing across soils, crops, and locations before routine application \parencite{Skaggs1995}, and using process-based models to feed databases for economic and risk analysis \parencite{Malone2007} are all advanced. Broader model-development needs include intense field measurement for validation and long-term climate-change simulation \parencite{Negm2014}, improved prediction of high-load peak events \parencite{Grewal2025}, DRAINMOD--DSSAT coupling to capture factors influencing maize relative yield under subirrigation \parencite{Gunn2018}, greater sensitivity to low flows and macropore-flow algorithms \parencite{Morrison2014}, incorporation of CO$_2$, disease, cultivar, and planting-date effects \parencite{Ndulue2024}, additional crop-genetic-parameter calibration and elevated-CO$_2$ accounting for winter crops \parencite{Li2025}, and improved subsurface-drainage simulation within soybean in APSIM \parencite{Dietzel2016}.

\subsubsection{Groundwater interactions and integrated transport}

Coupling drainage models to groundwater is a distinct need. \textcite{Addab2022} recommend SWAT-MODFLOW-RT3D for physically based flow and reactive transport once a salinity module is incorporated, and \textcite{Hoang2017} argue that improving SWAT's groundwater module to simulate water-table fluctuation and groundwater--surface-water exchange is necessary to represent nitrate retention in riparian zones realistically. Optimal regulation requires considering both spatial variation in N leaching and variation in N reduction across groundwater and surface-water systems \parencite{Hojberg2017}. Integrated mathematical models for surface-water--groundwater nutrient transport, with machine-learning application and climate-change evaluation \parencite{Elsayed2023}, defensible site-specific water balances to estimate the split between artificial drainage and groundwater flows \parencite{Barkle2021}, and improved NH$_4$-N transport and transformation modelling \parencite{Yang2008} are further priorities. Specialised reactive-transport extensions are also sought---for lime precipitation/dissolution and trace-element reactivity in salinity- and trace-element-affected systems \parencite{KARAJEH1994c}, and clay-mineralogy incorporation with field validation \parencite{Dadfar2010}.

\subsection{Greenhouse gases, emerging contaminants, and multi-objective tradeoffs}

An expanding frontier integrates greenhouse-gas (GHG) and contaminant considerations into drainage-mitigation assessment, especially the coupling between nitrate leaching and indirect nitrous oxide (N$_2$O). \textcite{Sawamoto2003} call for studies of the nitrate--N$_2$O relationship in groundwater and of indirect N$_2$O from subsurface drainage, and \textcite{Minamikawa2015} argue that new calculation methods for indirect N$_2$O from leaching and runoff are needed, along with evaluation of the tradeoff between indirect N$_2$O and N removal by heterotrophic denitrification across groundwater, drainage, river, and estuary \parencite{Minamikawa2013,Minamikawa2010}. Continued investigation of nitrate--N$_2$O linkages when N application and precipitation are temporally disconnected \parencite{OBrien2022}, the magnitude and fate of N$_2$O losses \parencite{Burton2024}, whether dissolved N$_2$O losses exceed surface emissions \parencite{Smith2019}, and long-term cumulative dissolved-GHG effects of management \parencite{Oliveira2025} are all open. \textcite{Webb2021} outline an artificial-water emissions-accounting agenda---seasonal emission-factor impacts, process-based modelling, gas-transfer-velocity models, the DOC/NO$_3$ ratio, and a global database---and \textcite{Fontes2019} recommend year-round monitoring of both N$_2$O and nitrate leaching for better total-N-loss estimates. Net-GHG evaluation in longer experiments \parencite{PrezaFontes2023}, the fate of waterborne methane and inclusion of GHGs alongside P balances when selecting measures \parencite{Carstensen2021}, and comparative life-cycle assessment based on complete N balances \parencite{HernandezRamirez2011} extend this frontier. In cold and poorly drained soils, biogeochemical measurement and modelling of methane production and oxidation is likewise needed \parencite{Kumar2014b}.

Emerging contaminants and biological interactions are also flagged. Better constraint of the mechanism and lifetime of glyphosate in soil and water \parencite{Spiese2023}, investigation of planter modifications to reduce aquatic exposure to neonicotinoid and other seed-applied pesticides \parencite{Schaafsma2019}, understanding of how watershed microbial communities respond to and control nutrient fluxes \parencite{Knorr2023}, and study of water-quality changes during different rainfall-runoff events to identify compound sources \parencite{Fraindová2022} are all advanced. Particulate silicon transport and its coupling to warming, phytolith mineralisation, and diatom dynamics \parencite{Sethna2022}, nutrient stoichiometry and the fate of organic matter in aquatic ecosystems \parencite{Manninen2023,Manninen2018}, and the release of accumulated compounds and antibiotic-resistant genes in sludge-amended soils \parencite{Ghirardini2019} broaden the constituent scope. Overarching tradeoffs remain unresolved: the tradeoff of increased nitrate losses under some practices \parencite{Derby2024}, the balance between crop-yield gains in tiled areas and downstream water-quality losses \parencite{Addab2022}, and unresolved tradeoffs among N-loss pathways when switching to fertigation \parencite{Gamble2018}. Emission- and load-optimisation is expected to hinge on optimal N rate and timing that minimise yield--environment tradeoffs \parencite{Emmett2022}.

\subsection{Monitoring, tracers, isotopes, and source apportionment}

Improved measurement and tracing underpin many proposed advances. \textcite{Williams2022} call for higher-frequency (sub-daily to daily) measurements, further work on groundwater dynamics in DRP transport, and more comprehensive source-apportionment methods, and \textcite{Hu2024} recommend quantifying $\delta^{18}$O of nitrification-derived nitrate across environments and coupling water and nitrate oxygen isotopes to quantify evaporative enrichment. Monitoring spatial heterogeneity and temporal variability of nitrate concentrations and isotopic compositions in topsoil, drains, and ditches, using multiple tracers, is deemed essential to characterise soil's contribution to river nitrate \parencite{Briand2017}, and boron-isotope analyses around feedlots and fields are advanced to strengthen contaminant tracing \parencite{Komor1997}. Combining isotopic and pathway-partitioning approaches---explicit reporting of surface-inlet presence and partitioning of water and nutrient loads by pathway via before-after control-impact designs, hydrograph separation, or tracers \parencite{Flores2021}, evaluating overland-flow versus tile-flow contributions to edge-of-field losses \parencite{Kokulan2019}, and stable water isotopes to verify groundwater flow paths and identify plant water sources \parencite{Islam2022,YliHalla2022}---is repeatedly urged. The role of groundwater outflow in nutrient loading is flagged as under-quantified: modelling and measurement of its contribution \parencite{Turunen2013}, its inclusion in future N-loss studies with attention to non-mineral fractions \parencite{Jokinen2024}, and accounting for the redox status of underlying groundwater when estimating total-N export \parencite{Barkle2021} are all needed. Finer temporal analyses---interactions between drain flow and concentration variation \parencite{Lagzdins2015}, whether near-term flow-weighted nitrate changes predict long-term load and yield changes \parencite{Danalatos2022}, and detailed storm-hydrograph analysis to test event-scale hypotheses \parencite{Kladivko2004}---complete this methodological frontier.

\subsection{Soil health, biology, and drainage-induced soil evolution}

Slower-acting soil and biological processes are repeatedly flagged as understudied. The soil-biology response to changing drainage hydrology requires longer-term evaluation of biological indicators under standardised sampling \parencite{Sherbine2023}, study of the interaction mechanisms between subsurface drainage and soil physicochemical properties and microbial activity \parencite{Yang2025}, and continued research on microbial community structure and function to explain tillage and drainage impacts \parencite{Kaur2024b,Kumar2014}. Scalable nematode measurement integrated with a soil-health index \parencite{Li2023b}, valid and repeatable soil-quality indicators within a systems framework \parencite{Nakajima2015}, soil-enzyme activities as management diagnostics \parencite{Singh2007}, proof that no-till improves soil-health functions enough to reduce nitrate leaching sustainably \parencite{Masri2024}, and dynamic compaction indicators with long-term forage or agroforestry restoration trials \parencite{Caron2025} are all advanced. Specific biological questions include the link between \textit{Fusarium} abundance and soybean root rot under long-term drainage \parencite{Han2017}, disruption of crack-to-burrow linkages through conservation tillage \parencite{Shipitalo2004}, precise evapotranspiration quantification to explain drain-flow differences \parencite{Daigh2014}, drought tolerance and phenology of rice cultivars guiding drainage-system selection \parencite{DarziNaftchali2017b}, and forage-sorghum nitrate testing before grazing \parencite{Nash2017}.

A distinctive line of inquiry treats subsurface drainage as a driver of long-term, potentially chaotic soil evolution. \textcite{Montagne2009} argue that quantifying human-induced soil evolution as a function of time should be a research priority, including drainage's impact on clay-mineral dissolution, and that soil evolution in response to climate change must be explicitly studied and incorporated in global models \parencite{Montagne2010}; significant material losses over a few decades make human-induced perturbation a powerful tool for studying soil-evolution timescales \parencite{Montagne2013}. Whether drainage-induced evolution continues or stabilises over decades, with measurements along drains at different slope positions and network ages, remains open \parencite{Frison2009}, and a long classical experiment continues to test whether soil structure is eventually affected \parencite{HOOGHOUDT1952}. Empirical prediction of earthworm activity from tillage, residue cover, and particle size---likely stronger for no-till than conventional systems---is proposed \parencite{Ouellet2008}, and maintaining drainage operation and cultivation is framed as a management imperative for long-term soil health \parencite{Singh2007}.

\subsection{Salinity, acid sulfate, reclamation, and irrigation-return management}

In arid, saline, and irrigated settings, distinct research and management needs emerge. \textcite{Sharma2004} suggest sequential reuse across cereals, grasses, forest plantations, and halophytes and innovative evaporation-pond designs, while gypsum, elemental sulphur, and lime amendments are recommended for long-term management of problem soils \parencite{MANN1982,He2015}. Tolerant rootstocks, halophyte germplasm, and transgenic crops \parencite{Cuevas2019}, land-treatment/pond combinations with restricted hydraulic loading \parencite{Christen2010}, pre-treatment and preferential-flow control before wastewater application \parencite{Duan2015}, and biochar selected for low contaminant content with urgent long-term impact assessment \parencite{Tian2023b} are all advanced, alongside optimum rates of organic soil conditioners to minimise adverse water-quality impacts \parencite{Zhang2015}. Modelling and management priorities include long-term geohydrological and crop data to validate simulations across irrigation/drainage systems \parencite{Yao2017}, systematic evaluation of alternative field-scale irrigation and drainage strategies under ambiguity with a statistical likelihood of achieving goals \parencite{Alzraiee2013}, land-use planning steering irrigation toward low-salinity aquifers \parencite{Bethune2008}, salt-tolerant crops in reclamation rotations \parencite{Qadir2001}, and quantification of rainfall and irrigation-efficiency gains relative to salt-removal capacity \parencite{Castanheira2010}. Water-productivity and technology questions recur: HYDRUS--EPIC evaluation of saline-water-irrigation sustainability \parencite{Feng2019}, water-table contributions to storage and evapotranspiration \parencite{Jia2017}, scalable data infrastructure, life-cycle cost assessment, and low-temperature crystallisation toward near-zero discharge \parencite{Qin2025}, sensor- and expert-informed scheduling \parencite{Abdelkrim2021}, complete economic evaluation of irrigation-technology shifts \parencite{DINAR1989}, and continued monitoring of yield and water-quality change to forecast groundwater suitability \parencite{RADSTAKE1988,Karegoudar2022,Qian2025,ElHassanin2022,Ayars2006}. Sediment- and return-flow reduction from furrow-irrigated fields, improved pond removal, and nitrogen-leaching-limiting irrigation and nutrient management are prioritised \parencite{Bjorneberg2020}.

Drained acid sulfate (AS) and boreal soils raise specialised questions: short-term hydro-meteorological effects on sulfide oxidation, ripening effects on hydraulic properties, and secondary-mineral formation should be incorporated into model applications \parencite{Salo2023}; the role of the N cycle in sulfide oxidation is unknown in boreal cultivated AS soils \parencite{Virtanen2017}; the peat-cover thickness required to hinder sulfidic-material oxidation needs detailed study \parencite{YliHalla2022}; and more detailed AS-soil maps, understanding of anthropogenic versus organic acidity, and cessation of drainage of potential AS soils are urged \parencite{Åström2005}. For reclaimed saline, waterlogged, and Vertisol systems, priorities include how subsurface drainage integrated with crop and residue management affects soil quality, carbon sequestration, nutrient budgets, and GHG fluxes \parencite{Mukhopadhyay2023,Mukhopadhyay2023b}, long-term effects of controlled drainage on salinity \parencite{ElGhannam2023}, salt-crust formation and dissolved-species transport \parencite{Colombani2016}, leaching effects on additional elements and organic modifiers \parencite{SoltaniGerdefaramarzi2020}, monitoring of elemental toxicities and waterlogging-stage-specific drainage scheduling \parencite{Qian2021}, long-term root-growth resistance and drainable-porosity determination \parencite{Bahçeci2018,Bahceci2008}, additive-effectiveness and land-replacement characterisation \parencite{Isa2023}, and further multi-location, multi-condition testing over longer periods \parencite{Sadhukhan2019}. Retired land over shallow water tables surrounded by irrigated fields raises unresolved ultimate-use questions \parencite{Purkey2001}.

\subsection{Remote sensing, precision tools, decision support, and data infrastructure}

Emerging sensing and data technologies are widely seen as enabling the next phase of research. \textcite{Khanal2020} outline a broad agenda: optimal conditions for detecting tile-drainage patterns and automatic extraction of drain lines; finer-resolution soil moisture from multi-source downscaling; remote sensing coupled with machine learning for soil properties and compaction; within-field nitrogen-stress mapping for variable-rate N; and semi-supervised and generative approaches to weed identification and site-specific residue management. \textcite{Cassman2025} propose a spatial framework combining climate and soil type with geo-referenced farmer-reported data, tailoring a Technology Extrapolation Domain framework and using remote sensing and on-farm databases to estimate field-specific benefits. \textcite{Reid2018} call for validating aggregate edge-of-field losses for Canadian conditions, quantifying edge-of-field-to-surface-water delivery ratios, and intensity-based P-loss indicators, and \textcite{Julien2023} envision a soil-wetness-index approach within a decision-support tool, noting drainage-season and threshold assumptions need testing across soils. Integrated water-and-nutrient management with advancing precision agriculture \parencite{Jung2024}, remote-sensing testing for environmental monitoring \parencite{Fontes2019}, and downscaled climate information with better forecasts and understanding of how farmers incorporate excess-precipitation experience into decisions \parencite{Morton2015} extend the tooling frontier.

Numerous studies frame the frontier in terms of improved indices, tools, and shared data. Development of a phosphorus index for German conditions incorporating subsurface drainage, peat groundwater transport, and stream-network connectivity, possibly combined with an N index \parencite{Buczko2007}, and improvement of nutrient-index schemes to include wetland and riparian-buffer removal capacity and catchment soil wetness \parencite{Pärn2012} are proposed. Regression-based tools for nutrient-credit trading via DWM can be extended to other regions with user-friendly interfaces \parencite{Negm2016}, and modifying the Soil Vulnerability Index to differentiate surface from subsurface artificial drainage and incorporate irrigation and rainfall intensity is advanced \parencite{Baffaut2020}. Stress s-curve refinement for aeration effects on yield \parencite{Wang2016}, regionally adjustable yield-loss coefficients \parencite{ROBINSON1992}, calibration and validation across multiple sub-watersheds \parencite{Singh2023}, coupling models with tools such as INVEST to assess ecosystem-service tradeoffs \parencite{McLellan2015}, and focusing conservation assessment on managing potholes that drain to ditches \parencite{Smith2008} round out the modelling-and-data infrastructure agenda. Underlying all of these is the data imperative: better, longer, and more spatially representative datasets---spanning water and solute fluxes, GHG emissions, root and biomass development, and management history---are the foundational requirement for validating models, resolving tradeoffs, and delivering site-specific, climate-resilient guidance \parencite{Ma2025,SKAGGS1994,Deng2011}.

\subsection{Climate change, extremes, and system resilience}

Climate adaptation frames many priorities. \textcite{Saha2024} stress continuing to identify practices and combinations that increase the resilience of agricultural systems to climate extremes across additional soils and longer periods, and \textcite{Malone2020} lay out a modelling agenda addressing additional periods, representation concentration pathways, management options, cultivars, cover-crop types, and soils. Within-county variation in fertiliser rate, crop-specific manure applications, and multiple global-climate-model forcings from CMIP5 or CMIP6, with attention to perennial-grass co-benefits, are advanced \parencite{Ferin2023}, and cropping systems designed to function under both drought and flood---with engineering solutions such as capturing excess floodwater---are prioritised \parencite{Dietzel2016}. Costs and benefits of adaptation to changing precipitation, irrigation-investment payback, and carbon-farming knowledge-sharing \parencite{PeltonenSainio2021}, study of extreme-event effects on cover-crop impacts for field-scale adaptation \parencite{Gupta2023b}, and comparison of N-fertiliser consequences under warming with best field-management identification \parencite{Haghnazari2020} extend the climate frontier. Testing and refinement across additional rainfall environments, soils, and rotations---to confirm interactions between subsurface drainage and planting pattern \parencite{Takeshima2023}, validate pipe-burial-depth effects on water productivity \parencite{Liu2025}, and characterise responses across soil and rainfall levels \parencite{Grewal2025,Negm2014,BENOIT1967}---is a recurring extension of this theme.

\subsection{Economics, incentives, policy, adoption, and the human dimension}

A consistent frontier is that biophysical understanding alone is insufficient; economic evaluation, incentive design, policy reform, and farmer engagement are decisive for adoption. \textcite{Liu2021} call for advanced case-by-case cost analyses and inter-agency action to promote 4R implementation, \textcite{Yuan2022} for more data on the cost-effectiveness and ecological benefits of conservation practices, and \textcite{McLellan2015} for incorporating practice costs into future work. Continued positive incentives such as CREP funding for wetlands, more educational programmes to increase adoption, and greater consideration of non-excludable benefits in multifunctional landscapes are recommended \parencite{Christianson2014}, and \textcite{David2015} argue that incentivising on-farm research, generating a collective sense of conservation need, and scaling individual efforts to meso- and macro-scale policy are essential. Cost-share and risk-insurance mechanisms \parencite{Jaynes2012,Karlen2005}, legislative reassessment and revised incentives to remove barriers to natural-water-retention measures \parencite{Kvítek2023}, national policy change to retain open ditch-and-field networks with targeted agri-environmental support \parencite{Marja2013}, and nutrient-trading programmes for shallow-drain adoption \parencite{Ghane2021b} are all advanced.

Adoption, capacity, and social dimensions are stressed across contexts. Detailed cost-benefit analyses, joint farmer--researcher experimentation, and capacity building are prioritised, with adoption depending on farmers' long-term cost-benefit perspective \parencite{Okwany2016}; creating know-how and farmer confidence before project execution is emphasised \parencite{Ghumman2011}; and key practice-use issues---where to apply BMPs, tool availability, how to bundle, and transferability to commercial farms---are identified \parencite{Ward2018}. Stakeholder willingness to adopt must inform proposed measures \parencite{Cuevas2019}, and research integrating ecological, pollution-related, and economic impacts of alternative land uses, with models for optimal policy design, is urged \parencite{Saikkonen2014}. In salinity and irrigation contexts, packages spanning reduced rice plantings, farmer compensation, flexible cost-recovery, and gradual subsidy liberalisation \parencite{Kotb2000}, and analysis showing that feasible drainage-water recycling requires markedly higher, less volatile irrigated yields and lower adoption costs \parencite{Sellars2022} are advanced. Economics of cover crops \parencite{Waring2020}, quantification of soil-health and P-loss impacts to guide subsidies \parencite{Roth2018}, long-term cumulative yield-and-profit assessment of drainage installation \parencite{Mourtzinis2021}, producer-participation and pilot-programme information for policy optimisation \parencite{Malone2023}, integrated biophysical--economic modelling at field scale \parencite{Ferin2023}, replicable studies in the Lower Mississippi Alluvial Valley for policy analysis \parencite{Kröger2012}, and landscape-zone management with decision tools accounting for soil effects on yield \parencite{Bakhsh2007b,Nelson2011,Kanwar1997} together define the economic and policy dimension.

\subsection{Synthesis: converging priorities for the field}

Read together, these frontiers underscore a coherent research programme rather than a scattering of isolated gaps. The field repeatedly asks for the same foundational investments---longer, standardised, multi-site records that span climatic variability and slow soil evolution; richer metrics for phosphorus saturation, colloidal transport, and legacy nutrients; isotopic and tracer methods to resolve pathways and apportion sources; and process-based models whose weakest links (winter and frozen-soil processes, soil C--N cycling, macropore flow, and groundwater coupling) are strengthened and validated against better data. On the management side, the consistent message is that no single practice suffices: in-field nitrogen and cropping-system change, drainage-water and water-table management, and edge-of-field treatment must be bundled, sequenced, and spatially targeted to hydrologically active critical source areas, with explicit attention to tradeoffs among nitrate reduction, phosphorus loss, greenhouse-gas emissions, yield, and economic risk. Crucially, the biophysical questions are inseparable from human ones---cost-effectiveness, incentive design, policy reform, and farmer adoption---so that the decisive advances will come from interdisciplinary work linking mechanistic understanding, robust tools and data infrastructure, and economically and socially viable pathways to adoption \parencite{Ma2025,SKAGGS1994,Ferin2023,Kvítek2023,Baulch2019}.